\begin{ouverture}[courbe]{6}{Les fonctions de référence}{Partie 2 · Fonctions}
{Affine · Carré · Cube · Inverse · Racine · Valeur absolue}
\objectif{Reconnaître les courbes et leurs domaines.}
\objectif{Comparer des images avec une fonction adaptée.}
\objectif{Résoudre des équations et inéquations de référence.}
\prerequis{Images, antécédents, ordre et racines.}
\end{ouverture}
\pagestyle{otMain}\essentiel\fontsize{11}{15}\selectfont
\essSection{1}{Un répertoire de fonctions}
\begin{definition}Une fonction affine est de la forme $f(x)=mx+p$. Elle est linéaire si $p=0$. Sa courbe est une droite, de pente $m$ et d’ordonnée à l’origine $p=f(0)$. Pour $a\ne b$, le taux d’accroissement $\dfrac{f(b)-f(a)}{b-a}$ vaut toujours $m$ : quand $x$ augmente de 1, $f(x)$ varie de $m$.\end{definition}
\begin{propriete}
\begin{tabularx}{\linewidth}{@{}lYY@{}}\toprule
Fonction & Domaine & Repère utile\\\midrule
$x\mapsto x^2$ & $\mathbb R$ & Positive ; $(-x)^2=x^2$ : courbe symétrique par rapport à l’axe des ordonnées.\\
$x\mapsto x^3$ & $\mathbb R$ & Signe de $x$ ; $(-x)^3=-x^3$.\\
$x\mapsto1/x$ & $\mathbb R\setminus\{0\}$ & Signe de $x$ ; ne vaut jamais 0.\\
$x\mapsto\sqrt x$ & $[0;+\infty[$ & Positive ou nulle ; carré égal à $x$.\\
$x\mapsto|x|$ & $\mathbb R$ & Distance à 0 ; courbe en V.\\\bottomrule
\end{tabularx}
\end{propriete}
\begin{savoirfaire}[Valeur absolue]Pour $x\ge0$, $|x|=x$ ; pour $x\le0$, $|x|=-x$. La fonction est décroissante avant 0 et croissante après 0. Son minimum vaut 0.\end{savoirfaire}
\begin{attention}$\sqrt{x^2}=|x|$, et non toujours $x$. L’inverse n’est pas défini en 0 ; une racine carrée réelle n’est pas définie pour une entrée négative.\end{attention}
\clearpage
\essSection{2}{Utiliser les courbes et les variations}
\begin{propriete}
La fonction carré est décroissante sur $]-\infty;0]$ et croissante sur $[0;+\infty[$. La fonction inverse est décroissante sur chacun des intervalles $]-\infty;0[$ et $]0;+\infty[$. Les fonctions cube et racine carrée sont croissantes sur leur domaine.
\end{propriete}
\begin{propriete}[Signe d’une fonction affine]
Si $m\ne0$, $f(x)=mx+p$ s’annule seulement en $x_0=-p/m$ ; $f(x)$ est du signe de $m$ pour $x>x_0$ et du signe contraire pour $x<x_0$. La droite est au-dessus de l’axe des abscisses exactement là où $f(x)>0$. Si $m=0$, $f$ garde le signe de $p$.
\begin{center}\renewcommand{\arraystretch}{1.25}\begin{tabular}{l|ccccc}\toprule
$x$ & $-\infty$ & & $-p/m$ & & $+\infty$\\\midrule
$mx+p$ avec $m>0$ (croissante) & & $-$ & $0$ & $+$ & \\
$mx+p$ avec $m<0$ (décroissante) & & $+$ & $0$ & $-$ & \\\bottomrule
\end{tabular}\end{center}
\end{propriete}
\begin{savoirfaire}[Comparer des images]
Identifier un intervalle sur lequel la fonction est monotone. Pour $0<a<b$, $a^2<b^2$, $\sqrt a<\sqrt b$ et $1/a>1/b$. Pour deux nombres négatifs, la comparaison des carrés se renverse.
\end{savoirfaire}
\begin{exemple}$-5<-2$ mais $25>4$. Pour $x\ge0$, $x^2-x=x(x-1)$ : la courbe du carré est sous la droite $y=x$ entre 0 et 1, et au-dessus pour $x>1$.\end{exemple}
\begin{savoirfaire}[Résoudre]
Pour $a\ge0$, $x^2\le a$ donne $-\sqrt a\le x\le\sqrt a$. Pour $b\ge0$, $\sqrt x\le b$ donne $0\le x\le b^2$. Avec un quotient, utiliser un tableau de signes plutôt que multiplier par une expression dont le signe est inconnu.
\end{savoirfaire}
\otIntertitre{Images mentales}
Retenir l’allure de chaque courbe, puis l’associer à son domaine, son signe et ses variations. Les planches d’exercices permettent de les tracer avec des échelles explicites.
\clearpage\rubrique{Découverte}
\begin{blocLivre}{ex-6-1}
\begin{exercice}[1]{Une fonction affine}
Pour $f(x)=-2x+3$, donner la pente, l’ordonnée à l’origine et l’image de 4. Tracer la droite à l’aide de deux points.
\label{fig-ex-6-1}
\figureGeo{reference-affine}
\end{exercice}
\end{blocLivre}
\begin{blocLivre}{ex-6-2}
\begin{exercice}[1]{Comparer des carrés}
Comparer sans calculatrice $(-7)^2$ et $(-3)^2$, puis $0{,}4^2$ et $0{,}7^2$.
\end{exercice}
\end{blocLivre}
\begin{blocLivre}{ex-6-3}
\begin{exercice}[1]{La fonction inverse}
Calculer les images de $-4$, $1/2$ et 5 par la fonction inverse. Le nombre 0 a-t-il un antécédent ?
\label{fig-ex-6-3}
\figureGeo{reference-inverse}
\end{exercice}
\end{blocLivre}
\begin{blocLivre}{ex-6-4}
\begin{exercice}[1]{Racine et domaine}
Résoudre $\sqrt x=5$, puis $\sqrt x=-2$.
\label{fig-ex-6-4}
\figureGeo{reference-racine}
\end{exercice}
\end{blocLivre}
\begin{blocLivre}{ex-6-5}
\begin{exercice}[1]{La courbe en V}
Calculer $|-3|$, $|0|$ et $|2{,}5|$. Résoudre $|x|\le4$.
\label{fig-ex-6-5}
\figureGeo{reference-absolue}
\end{exercice}
\end{blocLivre}

\clearpage\rubrique{Entraînement}
\begin{blocLivre}{ex-6-6}
\begin{exercice}[2]{Une équation de cube}
Résoudre $x^3=-27$, puis comparer $(-1{,}2)^3$ et $(-1{,}1)^3$.
\label{fig-ex-6-6}
\figureGeo{reference-cube}
\end{exercice}
\end{blocLivre}
\begin{blocLivre}{ex-6-7}
\begin{exercice}[2]{Une inéquation de carré}
Résoudre $x^2<9$ et $x^2\ge9$.
\label{fig-ex-6-7}
\figureGeo{reference-carre}
\end{exercice}
\end{blocLivre}
\begin{blocLivre}{ex-6-8}
\begin{exercice}[2]{Comparer des inverses}
Comparer $1/3$ et $1/7$, puis $1/(-5)$ et $1/(-2)$. Expliquer pourquoi on traite les deux paires séparément.
\end{exercice}
\end{blocLivre}
\begin{blocLivre}{ex-6-9}
\begin{exercice}[2]{Carré et identité}
Pour $x\ge0$, résoudre $x^2\le x$ et préciser les cas d’égalité.
\label{fig-ex-6-9}
\figureGeo{carre-identite}
\end{exercice}
\end{blocLivre}
\begin{blocLivre}{ex-6-10}
\begin{exercice}[2]{Une inéquation d’inverse}
Résoudre dans $\mathbb R$ l’inéquation $1/x\ge2$.
\end{exercice}
\end{blocLivre}

\clearpage\rubrique{Problèmes}
\begin{blocLivre}{ex-6-11}
\begin{exercice}[3]{Une distance minimale}
On considère $f(x)=|x-3|$ sur $[0;8]$. Déterminer son minimum et son maximum et préciser où ils sont atteints.
\end{exercice}
\end{blocLivre}
\begin{blocLivre}{ex-6-12}
\begin{exercice}[3]{Un problème de surface}
Une plaque carrée a une aire comprise entre 25 et 81 cm$^2$. Donner l’intervalle possible pour son côté, puis pour son périmètre.
\end{exercice}
\end{blocLivre}
\begin{blocLivre}{ex-6-13}
\begin{exercice}[3]{Une équation mixte}
Résoudre $\sqrt x=x$ dans $\mathbb R$. Justifier que les transformations n’introduisent pas de solution parasite.
\end{exercice}
\end{blocLivre}
\begin{blocLivre}{ex-6-14}
\begin{exercice}[3]{Une affirmation sur les fonctions}
« Si $a<b$, alors $a^2<b^2$ et $1/a>1/b$. » Pour chaque conclusion, donner un contre-exemple et un domaine sur lequel elle est correcte.
\end{exercice}
\end{blocLivre}
\begin{blocLivre}{ex-6-15}
\begin{exercice}[3]{Un encadrement transformé}
On sait $2\le x\le5$. Encadrer $x^2$, $\sqrt x$ et $1/x$. Les bornes obtenues sont-elles atteignables ?
\end{exercice}
\end{blocLivre}

\clearpage
